% GENERATED FILE. Edit canonical lessons, then run npm run generate:books.
% canonical-source-hash: fnv1a64:dea72404dd8b1daa
% canonical-lessons: JA-C485-oru, JA-C485-kinu, JA-C485-nunoji, JA-C485-jimi, JA-C485-hade, JA-R485-review-480-485

\chapter{Silk and Cloth}
\label{ch:ja-silk-and-cloth}
\hlchaptermodality{485}

\begin{chapteropening}
\textbf{By the end of this chapter:} \emph{I can say what cloth is like.}

Complete the last lesson of chapter 485: I can say what cloth is like.
\end{chapteropening}

\section[\texorpdfstring{\ja{おる} (oru)}{おる (oru)}]{\ja{おる} (oru) — to weave}

\label{lesson:JA-C485-oru}



\begin{quote}
Before the new one: say the Japanese for to come apart at the seams, then the Japanese for to tie up.
\end{quote}

\subsection*{You'll want to know: \ja{おる}}
\textbf{\ja{おる}} — \emph{oru} — \textquotedblleft{}to weave\textquotedblright{}.

\begin{quote}
\textbf{\ja{そぼは} \ja{きものの} \ja{ぬのを} \ja{おります。}} — \emph{sobo wa kimono no nuno o orimasu} — \textquotedblleft{}My grandmother weaves cloth for kimono.\textquotedblright{}
\end{quote}

Polite \textbf{\ja{おります}}, \emph{orimasu}; a godan verb. Past (polite): \textbf{\ja{そぼが} \ja{ぬのを} \ja{おりました。}}, \emph{sobo ga nuno o orimashita}.

\begin{cousinweb}
\textbf{\ja{おる}} (\emph{oru}) is a native Japanese word (wago), inherited rather than borrowed from Chinese; its kanji was attached later as a kun'yomi, and its older origin is not settled (JMdict, KANJIDIC2).
\end{cousinweb}

\begin{grammarlens}[title={What you've built}]
One more word for describing what you lived and felt.
\end{grammarlens}

\subsection*{Your turn}
\begin{itemize}
\raggedright
  \item \emph{Say it:} \emph{oru}
  \item \emph{Say it:} \emph{oru}, once more
  \item \emph{Say it:} \emph{sobo wa kimono no nuno o orimasu} — My grandmother weaves cloth for kimono.
  \item \emph{Recall:} say the Japanese for to come apart at the seams, then the Japanese for to tie up, then say \emph{oru} again
\end{itemize}

\begin{tcolorbox}[breakable,colback=teal!4,colframe=teal!35!black,title={Before you move on}]
What does \ja{おる} mean? (\textquotedblleft{}To weave\textquotedblright{}.) Say it once more.
\end{tcolorbox}
\section[\texorpdfstring{\ja{きぬ} (kinu)}{きぬ (kinu)}]{\ja{きぬ} (kinu) — silk}

\label{lesson:JA-C485-kinu}



\begin{quote}
Before the new one: say the Japanese for to tie up, then the Japanese for to weave.
\end{quote}

\subsection*{You'll want to know: \ja{きぬ}}
\textbf{\ja{きぬ}} — \emph{kinu} — \textquotedblleft{}silk\textquotedblright{}.

\begin{quote}
\textbf{\ja{きぬの} \ja{きものは} \ja{たかいです。}} — \emph{kinu no kimono wa takai desu} — \textquotedblleft{}Silk kimono are expensive.\textquotedblright{}
\end{quote}

\begin{cousinweb}
\textbf{\ja{きぬ}} (\emph{kinu}) is a native Japanese word (wago), inherited rather than borrowed from Chinese; its kanji was attached later as a kun'yomi, and its older origin is not settled (JMdict, KANJIDIC2).
\end{cousinweb}

\begin{grammarlens}[title={What you've built}]
One more word for describing what you lived and felt.
\end{grammarlens}

\subsection*{Your turn}
\begin{itemize}
\raggedright
  \item \emph{Say it:} \emph{kinu}
  \item \emph{Say it:} \emph{kinu}, once more
  \item \emph{Say it:} \emph{kinu no kimono wa takai desu} — Silk kimono are expensive.
  \item \emph{Recall:} say the Japanese for to tie up, then the Japanese for to weave, then say \emph{kinu} again
\end{itemize}

\begin{tcolorbox}[breakable,colback=teal!4,colframe=teal!35!black,title={Before you move on}]
What does \ja{きぬ} mean? (\textquotedblleft{}Silk\textquotedblright{}.) Say it once more.
\end{tcolorbox}
\section[\texorpdfstring{\ja{ぬのじ} (nunoji)}{ぬのじ (nunoji)}]{\ja{ぬのじ} (nunoji) — fabric; dress material}

\label{lesson:JA-C485-nunoji}



\begin{quote}
Before the new one: say the Japanese for to weave, then the Japanese for silk.
\end{quote}

\subsection*{You'll want to know: \ja{ぬのじ}}
\textbf{\ja{ぬのじ}} — \emph{nunoji} — \textquotedblleft{}fabric; dress material\textquotedblright{}.

\begin{quote}
\textbf{\ja{この} \ja{ぬのじで} \ja{ふくを} \ja{つくります。}} — \emph{kono nunoji de fuku o tsukurimasu} — \textquotedblleft{}I make clothes from this fabric.\textquotedblright{}
\end{quote}

\begin{cousinweb}
\textbf{\ja{ぬのじ}} (\emph{nunoji}) mixes the two layers, native and Sino-Japanese: native \emph{nuno}, \textquotedblleft{}linen\textquotedblright{} and Sino-Japanese \emph{ji}, \textquotedblleft{}ground\textquotedblright{} (Middle Chinese \emph{dijH}), borrowed from Chinese (JMdict, KANJIDIC2; Middle Chinese in Baxter's transcription).
\end{cousinweb}

\begin{grammarlens}[title={What you've built}]
One more word for describing what you lived and felt.
\end{grammarlens}

\subsection*{Your turn}
\begin{itemize}
\raggedright
  \item \emph{Say it:} \emph{nunoji}
  \item \emph{Say it:} \emph{nunoji}, once more
  \item \emph{Say it:} \emph{kono nunoji de fuku o tsukurimasu} — I make clothes from this fabric.
  \item \emph{Recall:} say the Japanese for to weave, then the Japanese for silk, then say \emph{nunoji} again
\end{itemize}

\begin{tcolorbox}[breakable,colback=teal!4,colframe=teal!35!black,title={Before you move on}]
What does \ja{ぬのじ} mean? (\textquotedblleft{}Fabric; dress material\textquotedblright{}.) Say it once more.
\end{tcolorbox}
\section[\texorpdfstring{\ja{じみ} (jimi)}{じみ (jimi)}]{\ja{じみ} (jimi) — plain; quiet (in colour)}

\label{lesson:JA-C485-jimi}



\begin{quote}
Before the new one: say the Japanese for silk, then the Japanese for fabric.
\end{quote}

\subsection*{You'll want to know: \ja{じみ}}
\textbf{\ja{じみ}} — \emph{jimi} — \textquotedblleft{}plain; quiet (in colour)\textquotedblright{}.

\begin{quote}
\textbf{\ja{ちちの} \ja{ふくは} \ja{いつも} \ja{じみです。}} — \emph{chichi no fuku wa itsumo jimi desu} — \textquotedblleft{}My father's clothes are always plain.\textquotedblright{}
\end{quote}

A na-adjective: \ja{な} before a noun.

\begin{cousinweb}
\textbf{\ja{じみ}} (\emph{jimi}) is Sino-Japanese (kango): two Chinese characters read with their on'yomi, readings borrowed from Chinese. They are \emph{ji}, \textquotedblleft{}ground\textquotedblright{} (Middle Chinese \emph{dijH}) and \emph{mi}, \textquotedblleft{}flavor\textquotedblright{} (\emph{mj+jH}) (KANJIDIC2; Middle Chinese in Baxter's transcription).
\end{cousinweb}

\begin{grammarlens}[title={What you've built}]
One more word for describing what you lived and felt.
\end{grammarlens}

\subsection*{Your turn}
\begin{itemize}
\raggedright
  \item \emph{Say it:} \emph{jimi}
  \item \emph{Say it:} \emph{jimi}, once more
  \item \emph{Say it:} \emph{chichi no fuku wa itsumo jimi desu} — My father's clothes are always plain.
  \item \emph{Recall:} say the Japanese for silk, then the Japanese for fabric, then say \emph{jimi} again
\end{itemize}

\begin{tcolorbox}[breakable,colback=teal!4,colframe=teal!35!black,title={Before you move on}]
What does \ja{じみ} mean? (\textquotedblleft{}Plain; quiet (in colour)\textquotedblright{}.) Say it once more.
\end{tcolorbox}
\section[\texorpdfstring{\ja{はで} (hade)}{はで (hade)}]{\ja{はで} (hade) — showy; flashy}

\label{lesson:JA-C485-hade}



\begin{quote}
Before the new one: say the Japanese for fabric, then the Japanese for plain.
\end{quote}

\subsection*{You'll want to know: \ja{はで}}
\textbf{\ja{はで}} — \emph{hade} — \textquotedblleft{}showy; flashy\textquotedblright{}.

\begin{quote}
\textbf{\ja{あの} \ja{ゆかたは} \ja{はでですね。}} — \emph{ano yukata wa hade desu ne} — \textquotedblleft{}That yukata is really flashy.\textquotedblright{}
\end{quote}

A na-adjective: \ja{な} before a noun.

\begin{cousinweb}
\textbf{\ja{はで}} (\emph{hade}) mixes the two layers, native and Sino-Japanese: Sino-Japanese \emph{ha}, \textquotedblleft{}faction\textquotedblright{} (Middle Chinese \emph{pheaH}), borrowed from Chinese and native \emph{de}, \textquotedblleft{}hand\textquotedblright{} (JMdict, KANJIDIC2; Middle Chinese in Baxter's transcription).
\end{cousinweb}

\begin{grammarlens}[title={What you've built}]
One more word for describing what you lived and felt.
\end{grammarlens}

\subsection*{Your turn}
\begin{itemize}
\raggedright
  \item \emph{Say it:} \emph{hade}
  \item \emph{Say it:} \emph{hade}, once more
  \item \emph{Say it:} \emph{ano yukata wa hade desu ne} — That yukata is really flashy.
  \item \emph{Recall:} say the Japanese for fabric, then the Japanese for plain, then say \emph{hade} again
\end{itemize}

\begin{tcolorbox}[breakable,colback=teal!4,colframe=teal!35!black,title={Before you move on}]
What does \ja{はで} mean? (\textquotedblleft{}Showy; flashy\textquotedblright{}.) Say it once more.
\end{tcolorbox}
\section[(dialogue)]{Looking back: from Making Things to Silk and Cloth}

\label{lesson:JA-R485-review-480-485}



\begin{quote}
Say the Japanese for to make and an apprentice.
\end{quote}

\subsection*{Your turn}
Say these aloud:

\begin{itemize}
\raggedright
  \item to make, an apprentice, a master (of a craft), time and effort, to finish off
  \item skill, clumsy, technology, wisdom, to devise
  \item a nail, a saw, wire, made of wood, to stick into
  \item clay, china, a bowl, to harden, to set hard
  \item sewing, to knit, knitting wool, to come apart at the seams, to tie up
  \item to weave, silk, fabric, plain, showy
\end{itemize}

\begin{tcolorbox}[breakable,colback=teal!4,colframe=teal!35!black,title={Before you move on}]
To make? (\textbf{\ja{つくる}}, \emph{tsukuru}.) An apprentice? (\textbf{\ja{でし}}, \emph{deshi}.) A master (of a craft)? (\textbf{\ja{ししょう}}, \emph{shishō}.) Time and effort? (\textbf{\ja{てま}}, \emph{tema}.) To finish off? (\textbf{\ja{しあげる}}, \emph{shiageru}.) Skill? (\textbf{\ja{うでまえ}}, \emph{udemae}.) Clumsy? (\textbf{\ja{ぶきよう}}, \emph{bukiyō}.) Technology? (\textbf{\ja{ぎじゅつ}}, \emph{gijutsu}.) Wisdom? (\textbf{\ja{ちえ}}, \emph{chie}.) To devise? (\textbf{\ja{くふうする}}, \emph{kufū suru}.) A nail? (\textbf{\ja{くぎ}}, \emph{kugi}.) A saw? (\textbf{\ja{のこぎり}}, \emph{nokogiri}.) Wire? (\textbf{\ja{はりがね}}, \emph{harigane}.) Made of wood? (\textbf{\ja{もくぞう}}, \emph{mokuzō}.) To stick into? (\textbf{\ja{ささる}}, \emph{sasaru}.) Clay? (\textbf{\ja{ねんど}}, \emph{nendo}.) China? (\textbf{\ja{せともの}}, \emph{setomono}.) A bowl? (\textbf{\ja{うつわ}}, \emph{utsuwa}.) To harden? (\textbf{\ja{かためる}}, \emph{katameru}.) To set hard? (\textbf{\ja{かたまる}}, \emph{katamaru}.) Sewing? (\textbf{\ja{さいほう}}, \emph{saihō}.) To knit? (\textbf{\ja{あむ}}, \emph{amu}.) Knitting wool? (\textbf{\ja{けいと}}, \emph{keito}.) To come apart at the seams? (\textbf{\ja{ほころびる}}, \emph{hokorobiru}.) To tie up? (\textbf{\ja{しばる}}, \emph{shibaru}.) To weave? (\textbf{\ja{おる}}, \emph{oru}.) Silk? (\textbf{\ja{きぬ}}, \emph{kinu}.) Fabric? (\textbf{\ja{ぬのじ}}, \emph{nunoji}.) Plain? (\textbf{\ja{じみ}}, \emph{jimi}.) Showy? (\textbf{\ja{はで}}, \emph{hade}.)
\end{tcolorbox}
